% ═══════════════════════════════════════════════════════════════
% MUSTERLÖSUNG: Das Ohmsche Gesetz (Klasse 8)
% Lösungen in ROT
% ═══════════════════════════════════════════════════════════════

\documentclass[11pt, a4paper]{article}

% ═══════════════════════════════════════════════════════════════
% PAKETE
% ═══════════════════════════════════════════════════════════════

\usepackage[utf8]{inputenc}
\usepackage[T1]{fontenc}
\usepackage[ngerman]{babel}
\usepackage{helvet}
\renewcommand{\familydefault}{\sfdefault}

\usepackage[margin=2cm, top=2.5cm, bottom=2cm]{geometry}
\usepackage{amsmath, amssymb}
\usepackage{siunitx}
\sisetup{locale = DE, per-mode = symbol}

\usepackage{graphicx}
\usepackage{tikz}
\usetikzlibrary{circuits.ee.IEC, positioning}
\usepackage{enumitem}
\usepackage{tabularx}
\usepackage{booktabs}
\usepackage{array}
\usepackage{multicol}

\usepackage{xcolor}
\usepackage{tcolorbox}
\usepackage{fancyhdr}
\usepackage{lastpage}

% Kompaktere Listen
\setlist{nosep, leftmargin=1.5em}

% ═══════════════════════════════════════════════════════════════
% FARBEN
% ═══════════════════════════════════════════════════════════════

\definecolor{physik}{HTML}{2c5aa0}
\definecolor{grau}{HTML}{888888}
\definecolor{hellgrau}{HTML}{f0f0f0}
\definecolor{success}{HTML}{2a7a4b}
\definecolor{stern1}{HTML}{2a7a4b}
\definecolor{stern2}{HTML}{b35c00}
\definecolor{stern3}{HTML}{c41e3a}
\definecolor{loesung}{HTML}{c62828}  % Rot für Lösungen

% ═══════════════════════════════════════════════════════════════
% VARIABLEN
% ═══════════════════════════════════════════════════════════════

\newcommand{\thema}{Das Ohmsche Gesetz}
\newcommand{\fach}{Physik}
\newcommand{\klasse}{8}
\newcommand{\stunde}{01}

% ═══════════════════════════════════════════════════════════════
% KOPF- UND FUSSZEILE
% ═══════════════════════════════════════════════════════════════

\pagestyle{fancy}
\fancyhf{}
\renewcommand{\headrulewidth}{0.4pt}
\renewcommand{\footrulewidth}{0pt}

\lhead{\small\textcolor{physik}{\textbf{\fach}} Kl.\,\klasse{} -- \thema{} \textcolor{loesung}{\textbf{-- MUSTERLÖSUNG}}}
\rhead{\small Name: \rule{5cm}{0.4pt}}

\cfoot{\small\thepage/\pageref{LastPage}}

% ═══════════════════════════════════════════════════════════════
% BEFEHLE
% ═══════════════════════════════════════════════════════════════

% Lösung (rot, unterstrichen)
\newcommand{\lsg}[1]{\textcolor{loesung}{\textbf{\underline{#1}}}}

% Formelbox
\newtcolorbox{formelbox}{
    colback=hellgrau,
    colframe=physik,
    boxrule=1pt,
    arc=0pt,
    left=8pt, right=8pt, top=8pt, bottom=8pt
}

% Merksatz
\newtcolorbox{merksatz}{
    colback=hellgrau,
    colframe=success,
    boxrule=1pt,
    arc=0pt,
    left=8pt, right=8pt, top=6pt, bottom=6pt,
    fonttitle=\bfseries,
    title=Merksatz
}

% Niveau-Marker
\newcommand{\niveauEins}{\textcolor{stern1}{\textbf{(*)}}}
\newcommand{\niveauZwei}{\textcolor{stern2}{\textbf{(**)}}}
\newcommand{\niveauDrei}{\textcolor{stern3}{\textbf{(***)}}}}

% Punkteangabe
\newcommand{\punkte}[1]{\hfill\textcolor{grau}{\small[#1 BE]}}

% ═══════════════════════════════════════════════════════════════
% DOKUMENT
% ═══════════════════════════════════════════════════════════════

\begin{document}

% ═══════════════════════════════════════════════════════════════
% KOPFBEREICH
% ═══════════════════════════════════════════════════════════════

\begin{center}
    {\Large\textbf{\textcolor{physik}{Arbeitsblatt: \thema}}}
    \quad
    {\Large\textcolor{loesung}{\textbf{-- MUSTERLÖSUNG}}}
\end{center}

\vspace{0.5em}

\begin{tcolorbox}[colback=hellgrau, colframe=grau, boxrule=0.5pt, arc=0pt]
\small
\textbf{Niveaus:}
\niveauEins{} Basis \quad
\niveauZwei{} Standard \quad
\niveauDrei{} Erweitert
\hfill
\textbf{Gesamtpunktzahl:} \lsg{25} / 25 BE
\end{tcolorbox}

\vspace{1em}

% ═══════════════════════════════════════════════════════════════
% TEIL 1: MERKSATZ
% ═══════════════════════════════════════════════════════════════

\section*{1. Formel und Merksatz}

\begin{formelbox}
\textbf{Das Ohmsche Gesetz:}

\vspace{1em}
\begin{center}
{\Large $U = R \cdot I$}
\end{center}

\vspace{0.5em}

\begin{tabular}{lll}
$U$ & = & Spannung in \lsg{Volt} (Einheit: \lsg{V}) \\[0.8em]
$R$ & = & \lsg{Widerstand} in Ohm (Einheit: \lsg{$\Omega$}) \\[0.8em]
$I$ & = & Stromstärke in \lsg{Ampere} (Einheit: \lsg{A}) \\
\end{tabular}
\end{formelbox}

\vspace{0.5em}

\begin{merksatz}
Je \lsg{höher} die Spannung, desto \lsg{höher} die Stromstärke.

Der Widerstand $R$ bleibt dabei \lsg{konstant} (bei ohmschen Widerständen).
\end{merksatz}

\vspace{1em}

% ═══════════════════════════════════════════════════════════════
% TEIL 2: MESSWERTE
% ═══════════════════════════════════════════════════════════════

\section*{2. Messwerte aus dem Experiment \niveauEins}

\textbf{Widerstand:} $R = \SI{20}{\ohm}$

\vspace{0.5em}

\begin{center}
\begin{tabular}{|c|c|c|c|c|c|}
\hline
\textbf{Spannung $U$} & \SI{2}{\volt} & \SI{4}{\volt} & \SI{6}{\volt} & \SI{8}{\volt} & \SI{10}{\volt} \\
\hline
\textbf{Stromstärke $I$} & \lsg{\SI{0,1}{\ampere}} & \lsg{\SI{0,2}{\ampere}} & \lsg{\SI{0,3}{\ampere}} & \lsg{\SI{0,4}{\ampere}} & \lsg{\SI{0,5}{\ampere}} \\
\hline
\end{tabular}
\end{center}

\textcolor{loesung}{\small Berechnung: $I = \frac{U}{R} = \frac{U}{\SI{20}{\ohm}}$, z.B. $I = \frac{\SI{2}{\volt}}{\SI{20}{\ohm}} = \SI{0,1}{\ampere}$}

\vspace{1em}

% ═══════════════════════════════════════════════════════════════
% TEIL 3: AUFGABEN NIVEAU *
% ═══════════════════════════════════════════════════════════════

\section*{3. Aufgaben}

\subsection*{Niveau \niveauEins{} -- Basis}

\textbf{Aufgabe 1:} Ergänze den Lückentext. \punkte{3}

Wenn man die Spannung am Widerstand \textbf{verdoppelt}, dann \lsg{verdoppelt} sich auch

die \lsg{Stromstärke}.

Man sagt: Die Stromstärke ist \lsg{proportional} zur Spannung.

\vspace{1.5em}

\textbf{Aufgabe 2:} Berechne die Spannung $U$. \punkte{2}

\begin{tcolorbox}[colback=white, colframe=physik!50, boxrule=0.5pt, arc=0pt]
Gegeben: $R = \SI{10}{\ohm}$, \quad $I = \SI{2}{\ampere}$

\vspace{0.5em}
Gesucht: $U = \,?$

\vspace{0.5em}
Formel: $U = R \cdot I$

\vspace{0.5em}
Einsetzen: $U = $ \lsg{\SI{10}{\ohm}} $\cdot$ \lsg{\SI{2}{\ampere}} $=$ \lsg{\SI{20}{\volt}}
\end{tcolorbox}

\vspace{1em}

% ═══════════════════════════════════════════════════════════════
% TEIL 4: AUFGABEN NIVEAU **
% ═══════════════════════════════════════════════════════════════

\subsection*{Niveau \niveauZwei{} -- Standard}

\textbf{Aufgabe 3:} Berechne die fehlenden Größen. \punkte{6}

\begin{enumerate}[label=\alph*)]
    \item $R = \SI{50}{\ohm}$, \quad $I = \SI{0,2}{\ampere}$ \quad $\Rightarrow$ \quad $U = $ \lsg{\SI{10}{\volt}}

    \textcolor{loesung}{\small $U = R \cdot I = \SI{50}{\ohm} \cdot \SI{0,2}{\ampere} = \SI{10}{\volt}$}

    \vspace{0.5em}

    \item $U = \SI{12}{\volt}$, \quad $R = \SI{4}{\ohm}$ \quad $\Rightarrow$ \quad $I = $ \lsg{\SI{3}{\ampere}}

    \textcolor{loesung}{\small $I = \frac{U}{R} = \frac{\SI{12}{\volt}}{\SI{4}{\ohm}} = \SI{3}{\ampere}$}

    \vspace{0.5em}

    \item $U = \SI{10}{\volt}$, \quad $I = \SI{0,5}{\ampere}$ \quad $\Rightarrow$ \quad $R = $ \lsg{\SI{20}{\ohm}}

    \textcolor{loesung}{\small $R = \frac{U}{I} = \frac{\SI{10}{\volt}}{\SI{0,5}{\ampere}} = \SI{20}{\ohm}$}
\end{enumerate}

\vspace{1em}

\textbf{Aufgabe 4:} Ordne zu: Welche Einheit gehört zu welcher Größe? \punkte{3}

\begin{center}
\begin{tabular}{ccc}
\textbf{Größe} & & \textbf{Einheit} \\[0.5em]
Spannung $U$ & \lsg{Volt (V)} & Ampere (A) \\
Stromstärke $I$ & \lsg{Ampere (A)} & Ohm ($\Omega$) \\
Widerstand $R$ & \lsg{Ohm ($\Omega$)} & Volt (V) \\
\end{tabular}
\end{center}

\vspace{1em}

% ═══════════════════════════════════════════════════════════════
% TEIL 5: AUFGABEN NIVEAU ***
% ═══════════════════════════════════════════════════════════════

\subsection*{Niveau \niveauDrei{} -- Erweitert}

\textbf{Aufgabe 5:} Fehlersuche \punkte{4}

Ein Schüler rechnet:

\begin{quote}
``$R = \SI{50}{\ohm}$, $I = \SI{200}{\milli\ampere}$.

Dann ist $U = 50 \cdot 200 = \SI{10000}{\volt}$.''
\end{quote}

a) Was hat der Schüler falsch gemacht?

\textcolor{loesung}{Der Schüler hat vergessen, Milliampere in Ampere umzurechnen.}

\textcolor{loesung}{$\SI{200}{\milli\ampere} = \SI{0,2}{\ampere}$ (nicht 200!)}

\vspace{0.5em}

b) Berechne das richtige Ergebnis.

\textcolor{loesung}{$I = \SI{200}{\milli\ampere} = \SI{0,2}{\ampere}$}

\textcolor{loesung}{$U = R \cdot I = \SI{50}{\ohm} \cdot \SI{0,2}{\ampere} = \SI{10}{\volt}$}

\vspace{1em}

\textbf{Aufgabe 6:} Transfer \punkte{4}

An einem Widerstand von $R = \SI{20}{\ohm}$ liegt eine Spannung von $U = \SI{6}{\volt}$ an.

\begin{enumerate}[label=\alph*)]
    \item Berechne die Stromstärke $I$.

    \textcolor{loesung}{$I = \frac{U}{R} = \frac{\SI{6}{\volt}}{\SI{20}{\ohm}} = \SI{0,3}{\ampere}$}

    \vspace{0.5em}

    \item Die Spannung wird auf $U = \SI{12}{\volt}$ verdoppelt.

    Was passiert mit der Stromstärke? Begründe ohne zu rechnen.

    \textcolor{loesung}{Die Stromstärke verdoppelt sich ebenfalls (auf $\SI{0,6}{\ampere}$),}

    \textcolor{loesung}{weil $I$ proportional zu $U$ ist (bei konstantem $R$).}
\end{enumerate}

\vspace{1em}

\textbf{Aufgabe 7:} Begründung \punkte{3}

Erkläre, warum der Widerstand $R$ konstant bleibt, auch wenn man die Spannung ändert.

\textcolor{loesung}{Der Widerstand ist eine Eigenschaft des Bauteils (Materialeigenschaft).}

\textcolor{loesung}{Er hängt von Material, Länge und Querschnitt des Leiters ab.}

\textcolor{loesung}{Die Spannung ist nur der ``Antrieb'' für den Strom, ändert aber nicht das Bauteil selbst.}

\vspace{1.5em}

% ═══════════════════════════════════════════════════════════════
% PUNKTEÜBERSICHT
% ═══════════════════════════════════════════════════════════════

\begin{tcolorbox}[colback=hellgrau, colframe=grau, boxrule=0.5pt, arc=0pt]
\small
\textbf{Punkteübersicht:}

\begin{tabular}{llcl}
\niveauEins{} Basis: & Aufgaben 1--2 & = & 5 BE \\
\niveauZwei{} Standard: & Aufgaben 3--4 & = & 9 BE \\
\niveauDrei{} Erweitert: & Aufgaben 5--7 & = & 11 BE \\
\hline
\textbf{Gesamt:} & & & \textbf{25 BE}
\end{tabular}
\hfill
\textbf{Erreicht:} \lsg{25} BE
\end{tcolorbox}

% ═══════════════════════════════════════════════════════════════
% BEWERTUNGSHINWEISE FÜR LEHRKRÄFTE
% ═══════════════════════════════════════════════════════════════

\newpage

\begin{center}
    {\Large\textbf{\textcolor{loesung}{Bewertungshinweise für Lehrkräfte}}}
\end{center}

\vspace{1em}

\subsection*{Punktevergabe}

\begin{tabular}{|l|c|l|}
\hline
\textbf{Aufgabe} & \textbf{BE} & \textbf{Kriterien} \\
\hline
1 (Lückentext) & 3 & Je 1 BE pro richtiger Lücke \\
\hline
2 (Einsetzen) & 2 & 1 BE Werte, 1 BE Ergebnis mit Einheit \\
\hline
3a (U berechnen) & 2 & 1 BE Rechnung, 1 BE Ergebnis mit Einheit \\
\hline
3b (I berechnen) & 2 & 1 BE Umstellung, 1 BE Ergebnis \\
\hline
3c (R berechnen) & 2 & 1 BE Umstellung, 1 BE Ergebnis \\
\hline
4 (Zuordnung) & 3 & Je 1 BE pro richtige Zuordnung \\
\hline
5a (Fehler finden) & 2 & Erkennen: mA nicht umgerechnet \\
\hline
5b (Korrektur) & 2 & 1 BE Umrechnung, 1 BE Ergebnis \\
\hline
6a (Berechnung) & 2 & 1 BE Formel/Rechnung, 1 BE Ergebnis \\
\hline
6b (Begründung) & 2 & Verdopplung + Proportionalität genannt \\
\hline
7 (Erklärung) & 3 & Materialeigenschaft, unabhängig von U \\
\hline
\end{tabular}

\vspace{1.5em}

\subsection*{Differenzierte Bewertung}

\textbf{Niveau (*) -- Berufsreife:}
\begin{itemize}
    \item Aufgaben 1--2 vollständig bearbeiten (5 BE)
    \item Ggf. Aufgabe 3a als Zusatz
    \item \textbf{Mindestanforderung:} 4 von 5 BE
\end{itemize}

\textbf{Niveau (**) -- Mittlere Reife:}
\begin{itemize}
    \item Aufgaben 1--4 bearbeiten (14 BE)
    \item \textbf{Mindestanforderung:} 10 von 14 BE
\end{itemize}

\textbf{Niveau (***) -- Gymnasium:}
\begin{itemize}
    \item Alle Aufgaben bearbeiten (25 BE)
    \item \textbf{Mindestanforderung:} 18 von 25 BE
\end{itemize}

\vspace{1.5em}

\subsection*{Häufige Fehler}

\begin{enumerate}
    \item \textbf{Einheiten vergessen:} Ergebnis ohne V, A, $\Omega$ $\Rightarrow$ halbe Punktzahl
    \item \textbf{mA nicht umgerechnet:} Typischer Fehler bei Aufgabe 5
    \item \textbf{Formel falsch umgestellt:} z.B. $I = U \cdot R$ statt $I = \frac{U}{R}$
    \item \textbf{Proportionalität nicht verstanden:} ``Strom bleibt gleich bei doppelter Spannung''
\end{enumerate}

\vspace{1.5em}

\subsection*{Nachteilsausgleich}

\textbf{LRS:}
\begin{itemize}
    \item Zeitzugabe: +10\%
    \item Rechtschreibfehler nicht werten
    \item Ggf. vergrößerte Kopie
\end{itemize}

\textbf{Förderbedarf Lernen:}
\begin{itemize}
    \item Nur Aufgaben 1--2 (Niveau *)
    \item Formelkarte als Hilfsmittel erlaubt
    \item Reduzierte BE-Anforderung
\end{itemize}

% ═══════════════════════════════════════════════════════════════
% ENDE
% ═══════════════════════════════════════════════════════════════

\end{document}
